\subsection{Exact structural and infinite-family construction filters}
The following results concern specified construction domains; they are not
general categorification exclusions for the rings that remain open.
The adjoint closure is the full ring for
$R8\_3\_65$, $R8\_3\_439$, $R8\_3\_567$,
$R8\_3\_901$ and $R8\_3\_1463$, excluding a nontrivial faithful group
grading of these particular categories. The two latter rings have a unique
Fibonacci generator. In a putative Fibonacci Deligne factorization, the
complementary simple basis would be forced to be $[0,1,2,7]$ and
$[0,1,3,6]$, respectively; neither basis is closed under multiplication.
They also cannot be braided: the Fibonacci subcategory would be
nondegenerate, since a symmetric fusion category is integral, and would
split off by \cite[Corollary~2.46 and Theorem~3.13]{DGNO}.
We do not infer a coherent central lift from based-ring commutativity.
The generalized near-group $R8\_3\_65$ has a free $C_4$ action on its four
noninvertibles and noninvertible FP dimension satisfying $d^2=1+11d$.
The classification \cite[Theorem~1]{SchopierayBraidedGNG} therefore excludes
a braided categorification as well.

There is a stronger unbraided construction restriction for $R8\_3\_65$.
Any canonical $\operatorname{Rep}(G)$ of an equivariantization must lie in
its pointed $C_4$ subcategory, since its other simple dimensions are
irrational. Hence every nontrivial possibility contains a central
$\operatorname{Rep}(C_2)$. Its sign object acts freely on all eight source
simples. The orbit description of equivariantization forces a rank-four
de-equivariantization of global dimension $125+55\sqrt5>24+12\sqrt3$,
contrary to the unrestricted rank-four bound proved in the rank-four
companion, Corollary~\texttt{finite:cor:sharp-bounds}. Thus every
nontrivial finite-group equivariantization is excluded for this ring.
In contrast, the complete rank-two based-module equations have exactly
two solutions up to exchanging vertices:
\[
N_\rho=\begin{pmatrix}0&1\\1&11\end{pmatrix},\qquad
N_\rho=\begin{pmatrix}3&5\\5&8\end{pmatrix}.
\]
An actual module category with either action would trivialize the pointed
$H^3(C_4,\mathbb C^\times)$ class. Neither that vanishing nor an
associator for the second action is inferred from these matrices.

For $R8\_3\_439$ and $R8\_3\_567$, the complete $C_3$ split-module
calculation permits exactly two conjugate nontrivial Fourier systems.
The subgroup algebra $A=\bigoplus_{g\in C_3}g$ exists because a fixed
noninvertible trivializes its pointed third-cohomology class. Its Morita
dual has six invertibles and two noninvertibles. The right deck actions
and the full left fusion action force the dual noninvertible matrix via
an integer linear system of rank $64$ in $64$ unknowns. The eight dual
actions are independent, and all $512$ resulting coefficients recover
the original respective source ring. This closes this particular Morita
route at the level of based rings; it asserts neither equivalence of
associators nor existence of a categorification.

We also enumerate every possible underlying object of a connected
separable algebra in $R8\_3\_65$, $439$, $455$, $567$, $901$ and $1463$.
The coefficient bounds $a_0=1$, $a_i=a_{i^*}$ and
$0\leq a_i\leq\lfloor\operatorname{FPdim}(X_i)\rfloor$, invertible
stabilization, and the entrywise nonnegative positive-semidefinite
induction residual Gram condition leave respectively
$30,28,21,22,18,26$ objects. Each entire finite domain is independently
re-enumerated with rational arithmetic; for the noncommutative tensors
the induction matrix uses right multiplication by the candidate algebra.
No surviving object is asserted to carry an associative multiplication.
For $R8\_3\_65$, every object $1+X$ with $X$ noninvertible fails this
necessary filter, leaving only the order-two pointed choice among
objects of the form $1+X$.

For the new cyclic family $\mathrm{NG}(C_n,n)$ of
\cite[Theorem~2.10 and Remark~2.11]{NearGroupTPE2026}, and indeed for every
categorification of those rings, the FP global dimension $D_n$ has
\[
\operatorname{Tr}(D_n)=n(n+4),\qquad
\operatorname{Nm}(D_n)=n^2(n+4).
\]
If a target quadratic FP dimension $D$ equals $sD_n$ with $s\in\mathbb Q_{>0}$,
then necessarily
\[
n=\frac{\operatorname{Tr}(D)^2}{\operatorname{Nm}(D)}-4,
\qquad
s=\frac{\operatorname{Tr}(D)}{n(n+4)}.
\]
Thus an exact calculation replaces a bounded search over $n$. Direct and
Morita equivalence require $s=1$; a group extension or equivariantization
requires integral $s$. If the target is a de-equivariantization by a group
of order $k=1/s$, then $k\mid n$, its rank is $n/k+k$, and its invertible
group is cyclic of order $n/k$. To see this, the integral subcategory
$\operatorname{Rep}(G)$ in the near-group category must be pointed, so
$G$ is cyclic. All its stabilizers have trivial second cohomology.
The equivariantization orbit description
\cite[Corollary~2.13 and Remark~3.1]{BurciuNataleEquivariantization}
then forces every invertible orbit to be fixed and the unique
noninvertible orbit to be free. Finally, $\mathrm{NG}(C_n,n)$ has trivial
universal grading, excluding proper neutral-component extraction.

These exact dimension, rank and invertible-group tests, supplemented by
the independent general exclusion of $R8\_3\_436$, exclude all five
single-operation routes for the original $34$ frontier tensors, for every
positive $n$. For example, $R8\_3\_65$ forces $(n,s)=(121,4/121)$;
$R8\_3\_439$ and $R8\_3\_567$ force $(n,k)=(12,2)$, whose
de-equivariantization would have cyclic, rather than $S_3$, invertibles.
Arbitrary chains of operations, Morita changes inserted into those chains,
and general noncentral algebra condensations are outside this claim.

For $R8\_3\_455$, a separate exact Fourier calculation excludes the
self-Morita-duality required of the neutral category within the strict
Evans--Gannon $\mathrm{HI}(C_3)$ matrix presentation, including its known
Galois companions. We retain the gap between an arbitrary abstract complex
neutral category with vanishing pointed cocycle obstructions and that
strict presentation. This conditional calculation does not change the
general status of $R8\_3\_455$.
